\section{Results}
\label{sec:results}

We report the four audit stages in the order of \S\ref{sec:method}.
Results use discovery data; aggregation, per-head estimates and sensitivity
analyses are detailed in Appendices~\ref{app:stage1}--\ref{app:stage4}.

\subsection{RI poorly identifies causal contributors}
\label{sec:results-why}

\paragraph{What RI selects.}
Under the descriptive percentile rule, the published pooled index selects two
heads whose qualifying attention is exclusively self- or previous-token
attention. Demonstrations contribute over 99\% of their summed scores
(Table~\ref{tab:pooled}), consistent with static token preferences rather
than evidence of retrieval. A separate matched-target test finds no head with
a significant preference after Holm correction; neither pooled-index head
has eligible matched events. The descriptive test-only rule selects
59 heads, with several candidates sensitive to the current token,
target anchor or fact order (Appendix~\ref{app:stage1}).

\begin{figure}[t]
\centering
\includegraphics[width=\columnwidth]{fig_ri_vs_importance}
\caption{RI versus signed causal importance (Scope~P, exact patching,
40 common pairs; symmetric-log axis). Orange: the combined 61-head
selection. (a) Pooled RI for all heads. (b) First-token name gap for
167 supported heads. Shading marks $|I_h|<0.3$.}
\label{fig:ri-vs-imp}
\end{figure}

\paragraph{Against the causal map.}
Pooled RI correlates weakly with signed importance ($\rho=0.13$),
while the test-only name gap correlates negatively
($\rho=-0.25$, and $-0.30$ within layer; Figure~\ref{fig:ri-vs-imp}).
The two pooled-index heads have near-zero measured effects.
Of the combined 61 selected heads, 48 have $|I_h|<0.1$,
less than 1\% of the clean--corrupted gap.
Among the 148 heads measured on all 178 pairs, 25 have
$|I_h|\geq0.3$, but only eight are selected.
The test-only selection also misses 88 of the 102 heads with the
largest positive common-subset effects.
Both selections are incomplete; the pooled selection additionally
meets our pre-registered misleading criterion
(Appendix~\ref{app:stage2}).

\paragraph{Why RI can miss retrieval.}
At the answer position, the current token is always a colon:
with unchanged visible token IDs, raw-embedding RI cannot adapt
its name preferences when mothers are reassigned.
All 49 jointly qualifying L16H21 clean/corrupted query pairs
exhibit this invariance.
A further hypothesis is that copying a current non-target name
raises the name-control score, penalizing potentially useful heads;
\S\ref{sec:results-char} examines this explanation.
These limitations concern the index, not the model's capacity
for semantic computation.

\subsection{The causal map}
\label{sec:results-map}

\paragraph{A few heads carry most of the effect.}
The gradient screen reproduces the exact ranking (Spearman 0.945 on the same
pairs), so the map below rests on exact patching wherever it is reported
(Appendix~\ref{app:stage2}). Against a mean clean--corrupted gap of 11.2 logits,
replacing a single head moves the metric by up to 5.7 logits (L17H1, half of the
gap), and three further heads each move it by about a quarter; all four have the
expected sign in every family (Figure~\ref{fig:causal-map}a). Twenty-five heads
reach $|I_h|\ge0.3$: 17 positive and 8 negative, the latter moving the model
\emph{toward} the stated fact when their corrupted-run output is inserted. Effects
of this size are not independent shares of the gap, since single-head patches can
interact. The same four heads lead the exact all-head scan on the common subset,
with effects within 5\% of their 178-pair values, so the leading heads are not an
artefact of the extension set (Appendix~\ref{app:stage2}).

\begin{figure}[t]
\centering
\includegraphics[width=\columnwidth]{fig_causal_map}
\caption{The causal map. (a) Signed Scope-P importance of the 25 strong heads on
all 178 pairs, with family-bootstrap intervals; orange marks RI-selected heads.
(b) Ratio of Scope-F to Scope-P importance by layer for the 20 heads of the
Scope-F set with $|I_h|>0.04$; marker size grows with $|I_h|$.}
\label{fig:causal-map}
\end{figure}

\paragraph{Two positions, two kinds of head.}
Restricting the patch to the final colon splits the map by depth
(Figure~\ref{fig:causal-map}b). Every measurable head in layers 16--26 keeps
between 91\% and 109\% of its effect when only the colon is patched, so its
contribution is delivered at the answer position; every measurable head in layers
6--11 loses all of it, so it acts only through what it writes at earlier test
positions for later components to read; two heads in layers 14--15 are
intermediate. The stored logits show how the
answer-position heads act: the positive ones raise the logit of the name retrieved
on the donor run and leave the alternative almost untouched, and the negative ones
lower it (Table~\ref{tab:scopeF-app}). The sign alone does not say by what
mechanism---promotion of an attended name, suppression of it, or an indirect
path---and the three strongest heads lie outside this Scope-F set;
\S\ref{sec:results-char} localizes every strong head's effect and tests these readings.

These positional differences suggest a writer--reader organization, but do not
identify computations or connections. The next stages characterize the heads
(\S\ref{sec:results-char}) and test causal routes between them
(\S\ref{sec:results-routes}).

\subsection{What the important heads do}
\label{sec:results-char}

\paragraph{Where the heads act, and what the writer writes.}
Twenty of the 25 strong heads deliver their whole effect at the colon
($I^F_h/I_h$ between 0.96 and 1.07; Appendix~\ref{app:stage3}). The exception is
the strongest head: L17H1 has no effect at the colon and acts inside the query
fact, at the mother's tokens, at ``is'' and at the period
(Figure~\ref{fig:profiles}a). At ``is'', ``of'' and the period the two prompts
have identical inputs, so the head's output there depends on the preceding mother
and matters for the answer. L15H25 acts at the last token of the query child. The
main writer thus sits in the same layer band as the readers, and the shallow heads
of \S\ref{sec:results-map} are all weak. A controlled readout of the residual at
``is'' exposes a mother-token preference that depends on L17H1: replacing the head
shifts the readout by $-0.46$ against $-0.52$ for the full source swap, negative in
all 20 families. This supports an identity-sensitive write, not a relational
binding; the content at L15H25's site is unresolved, since the readout does not
respond even to the source swap.

\begin{figure}[t]
\centering
\includegraphics[width=\columnwidth]{fig_head_profiles}
\caption{Head profiles. (a) Single-position patches of the three fact-position
heads, summed by token role in the query fact (M: mother tokens, C: child tokens,
oth.: the other test facts, Q: question, colon); family mean and SD over the 20
common families. (b) Attention mass at the colon of the strong answer-position
heads, split over the query mother, the query child, the three distractor mothers
and the colon itself, with each head's Scope-F importance.}
\label{fig:profiles}
\end{figure}

\paragraph{What the readers select and what they carry.}
The answer-position heads fall into three attention types
(Figure~\ref{fig:profiles}b): mother-attenders, of both signs; child-attenders,
which select the second mention of the queried child; and heads whose largest
target is the colon itself. For the mother- and child-attenders, changing only the
question moves attention to the newly queried fact in all 20 common families
(controls: at most a fifth of that shift), while reordering or swapping the mothers
leaves the family-mean attention almost unchanged. One negative head, L30H18,
attends to the distractor mothers. Separating attention from values at the colon,
value replacement carries most of the mean effect for the leading heads (L18H19:
values $+3.16$, pattern $+0.08$): the source selection follows the question, and
the swap acts through what is read there. Family-level interactions are large,
however, so separability is not assumed downstream. What the heads write is only
partly visible in their output vectors: L27H6 writes the answer directly ($+3.0$
logits for the query mother over the distractors), but for L18H19, L18H18, the
strong negative heads and the child-attenders the projection is far below the
effect. Their causal effects are not explained by this projection, which motivates
tests of downstream mediation. Recomputing the criterion from each head's actual
contextual input does not repair its correlation with importance ($-0.26\to-0.17$
on the 68 supported heads).

\paragraph{Copying is part of the story, not the mechanism.}
Six heads promote the first token of the name they attend to in at least 80\% of
events and one suppresses it; at the last-token anchor every label disappears. For
the positive movers the RI name gap is negative precisely at events whose current
token is a control name, as the copy-penalizing account of \S\ref{sec:results-why}
predicts. But copying is neither necessary nor sufficient (Appendix~\ref{app:stage3}):
the strongest weight copier hurts the answer by copying the distractor, and L23H15
favours the name it attends to while its effect is $-0.95$, so the sign of a head's
projection is not the sign of its effect. ``Negative mover'' covers at least two
patterns, and how L23H15 acts is open.

The profiles turn the map into testable connections: a fact-position writer whose
sites may feed the mother-attenders' keys or values, a second writer at the child
for the child-attenders, strong colon heads whose effects need a downstream path,
and positive and negative movers whose joint behaviour is untested
(\S\ref{sec:results-routes}).
